\documentclass[journal]{IEEEtran}

% IEEE starter checked on 2026-06-03 against:
% - IEEE Author Center guidance to use the IEEE Template Selector
% - CTAN IEEEtran package version 1.8b
% - CTAN bare_jrnl.tex skeleton
% Always replace this starter with the target journal's official current template
% before submission.

\usepackage{cite}
\usepackage{amsmath,amssymb,amsfonts}
\usepackage{algorithmic}
\usepackage{graphicx}
\usepackage{textcomp}
\usepackage{xcolor}
\usepackage{booktabs}
\usepackage{url}

\title{Journal Paper Title}

\author{First Author,~\IEEEmembership{Member,~IEEE,}
Second Author,~\IEEEmembership{Member,~IEEE}
\thanks{Manuscript received Month Day, Year; revised Month Day, Year.}
\thanks{The authors are with Department or Laboratory, University or Organization, City, Country (e-mail: email@example.com).}}

\begin{document}
\maketitle

\begin{abstract}
State the problem, the gap in existing work, the proposed approach, the key technical mechanism, the main quantitative result, and the implication.
\end{abstract}

\begin{IEEEkeywords}
IEEE, journal article, keyword one, keyword two
\end{IEEEkeywords}

\section{Introduction}
Introduce the technical problem, identify the gap, and summarize the contributions.

\section{Related Work}
Group prior work by technical theme and explain the distinction from this paper.

\section{Problem Formulation}
Define assumptions, variables, objectives, and constraints.

\section{Proposed Method}
Present the method, equations, algorithms, and implementation details needed for reproducibility.

\section{Experiments}
Describe scenarios or datasets, baselines, metrics, and implementation settings.

\section{Results and Discussion}
Report quantitative findings, connect them to the claims, and discuss limitations.

\section{Conclusion}
Summarize the technical contribution, main result, and implication.

\bibliographystyle{IEEEtran}
\bibliography{references}

\end{document}
